\chapter{Fuentes de datos y análisis preliminar}\label{chap:datos}

\section{Fuente principal: TLC Trip Record Data}

La fuente principal del trabajo es el portal de \textit{Trip Record Data} de la New York City Taxi and Limousine Commission. La página oficial ofrece descargas mensuales en formato Parquet y diccionarios de datos diferenciados por servicio~\cite{nyctlc_tripdata}. La TLC advierte además de que los datos se publican a partir de registros remitidos por operadores y bases, por lo que la publicación no debe interpretarse como garantía absoluta de exactitud de cada observación~\cite{nyctlc_reports}.

Esta característica refuerza la necesidad de una capa propia de calidad. El pipeline no modifica el fichero Bronze, pero registra incidencias de esquema y aplica reglas de validez antes de promover datos a Silver.

\section{Periodo de estudio}

La ventana general comprende desde el 1 de enero de 2018 hasta el 31 de diciembre de 2025. La elección permite dividir conceptualmente el análisis en tres fases:

\begin{itemize}
\item \textbf{Pre-pandemia:} 2018--2019, utilizado como referencia de comportamiento previo.
\item \textbf{Disrupción y transición:} 2020--2021, periodo en el que la movilidad experimenta cambios extraordinarios.
\item \textbf{Recuperación y nuevo equilibrio:} 2022--2025, donde se observa la evolución posterior y la consolidación de nuevas pautas.
\end{itemize}

La división es analítica, no causal. El trabajo no presupone que todo cambio entre periodos se deba a la pandemia.

\section{Disponibilidad por servicio}

La \cref{tab:services-coverage} resume el alcance previsto. La cobertura exacta se verificará automáticamente contra el portal durante la ejecución.

\begin{table}[H]
\centering
\caption{Servicios TLC incluidos y cobertura analítica prevista.}%
\label{tab:services-coverage}%
\begin{tabularx}{\textwidth}{lcccY}
\toprule
\textbf{Servicio} & \textbf{2018} & \textbf{2019--2025} & \textbf{Tarifas} & \textbf{Uso principal} \\
\midrule
Yellow & Sí & Sí & Amplias & Movilidad, economía y predicción \\
Green & Sí & Sí & Amplias & Movilidad, economía y predicción \\
FHV & Sí & Sí & Limitadas/variables & Movilidad y predicción \\
HVFHV & No separado & Sí & Esquema propio & Movilidad, comparación y predicción \\
\bottomrule
\end{tabularx}
\end{table}

Desde 2019 los HVFHV se reportan como un dataset separado y más detallado~\cite{nyctlc_tripdata}. En consecuencia, cualquier gráfico comparativo que incluya esta modalidad debe comenzar en 2019 o marcar explícitamente la falta de cobertura anterior.

\section{Unidad de ingesta}

La unidad de ingesta es el fichero mensual por servicio. Cada activo se identifica mediante una clave lógica compuesta por servicio, año, mes y URL de origen. Cuando sea viable se calcula además un hash criptográfico para detectar cambios en un fichero publicado con el mismo nombre.

Los metadatos mínimos son:

\begin{itemize}
\item servicio;
\item año y mes;
\item URL\@;
\item ruta local;
\item tamaño en bytes;
\item hash;
\item fecha de descubrimiento y descarga;
\item estado de validación;
\item estado de promoción a Silver;
\item número de registros leído y aceptado.
\end{itemize}

\section{Cambios de esquema}

Un histórico de ocho años no debe asumirse esquemáticamente estable. Los diccionarios TLC evolucionan y los campos disponibles dependen del servicio y periodo. Por ejemplo, la página oficial indica que a partir de 2025 se añadió un campo de cargo de congestión CBD en determinados datasets~\cite{nyctlc_tripdata}. Si el pipeline utilizara una lista rígida de columnas para todo el histórico, una incorporación de este tipo podría provocar errores o, peor aún, pérdidas silenciosas de información.

La estrategia adoptada separa tres conceptos:

\begin{enumerate}
\item \textbf{Esquema físico observado}: columnas y tipos contenidos en cada Parquet.
\item \textbf{Contrato Silver común}: conjunto de atributos homologados necesarios para análisis transversal.
\item \textbf{Atributos específicos}: campos propios de un servicio o periodo que se conservan cuando aportan valor.
\end{enumerate}

El contrato Silver permite que un consumidor conozca qué significa cada campo sin tener que estudiar cada versión del fichero original.

\section{Taxonomía común del viaje}

Se propone un núcleo de campos comunes como el mostrado en la \cref{tab:common-fields}. La existencia real de cada atributo se verificará por familia y año.

\begin{table}[H]
\centering
\caption{Contrato conceptual común para un viaje Silver.}%
\label{tab:common-fields}%
\begin{tabularx}{\textwidth}{l l Y}
\toprule
\textbf{Campo} & \textbf{Tipo} & \textbf{Significado} \\
\midrule
service\_type & texto & Modalidad TLC de procedencia. \\
pickup\_datetime & timestamp & Instante de inicio del viaje o recogida. \\
dropoff\_datetime & timestamp & Instante de finalización cuando esté disponible. \\
pickup\_location\_id & entero & Zona TLC de origen. \\
dropoff\_location\_id & entero & Zona TLC de destino. \\
trip\_distance & numérico & Distancia declarada cuando exista. \\
trip\_duration\_minutes & numérico & Duración derivada entre marcas temporales. \\
base\_fare & numérico & Tarifa base cuando la fuente la proporcione. \\
total\_amount & numérico & Importe total observado cuando sea comparable. \\
source\_file & texto & Fichero Bronze del que procede el registro. \\
source\_year/month & entero & Partición temporal de procedencia. \\
\bottomrule
\end{tabularx}
\end{table}

\section{Taxi Zone Lookup y geometrías}

La guía de usuario de la TLC proporciona una tabla de zonas y un fichero geográfico para representar las áreas utilizadas en los registros~\cite{nyctlc_userguide}. Se incorpora una dimensión \texttt{dim\_taxi\_zone} con identificador, nombre de zona, borough y \textit{service zone}. Las geometrías se conservan en un activo geográfico apropiado para construir mapas.

El uso de la taxonomía oficial evita inventar cuadrículas arbitrarias y facilita que los resultados puedan compararse con otros análisis basados en TLC zones.

\section{Calendario y festivos}

La dimensión de calendario se genera de forma determinista y no requiere un fichero masivo externo. Para cada fecha se incluyen, como mínimo, año, trimestre, mes, semana, día de la semana, indicador de fin de semana, estación y festivo. Los festivos se parametrizan con una fuente o librería versionada y deben quedar materializados en Gold para que un cambio de dependencia no modifique silenciosamente un experimento histórico.

\section{Meteorología}

Se prevé utilizar datos históricos de NOAA/NCEI para una estación representativa de Nueva York o una combinación documentada de estaciones~\cite{noaa_ncei}. La decisión final debe equilibrar disponibilidad temporal y simplicidad de unión. Para el objetivo predictivo, la granularidad preferente es horaria.

Las variables candidatas incluyen:

\begin{itemize}
\item temperatura;
\item precipitación;
\item nieve;
\item velocidad del viento;
\item visibilidad;
\item indicadores derivados de condiciones adversas.
\end{itemize}

Si una variable solo está disponible con granularidad diaria, se documentará para evitar dar a entender una precisión temporal inexistente.

\section{Reglas iniciales de calidad}

Antes de construir indicadores se aplican reglas generales y específicas. Entre las reglas generales se encuentran:

\begin{itemize}
\item marcas temporales dentro de rangos plausibles;
\item fecha de bajada posterior o igual a la de recogida cuando ambas existan;
\item duración no negativa;
\item distancias no negativas;
\item importes no negativos salvo campos donde un ajuste contable pueda justificarlo;
\item zonas de origen y destino pertenecientes al catálogo cuando el campo exista;
\item ausencia de duplicados según la clave técnica definida;
\item coherencia entre el mes del fichero y la fecha principal, permitiendo una tolerancia documentada en viajes que cruzan límites temporales.
\end{itemize}

Una regla que falla no implica siempre que el registro se elimine. Se distinguen errores críticos, advertencias y anomalías estadísticas. La causa de exclusión debe quedar registrada.

\section{Outliers y valores extremos}

Los datos de movilidad contienen valores extremos legítimos y errores. Un viaje largo hacia un aeropuerto o fuera de las zonas centrales no debe descartarse únicamente por ser infrecuente. Por esta razón se evita utilizar exclusivamente umbrales estadísticos globales para eliminar observaciones.

La estrategia combina restricciones físicas o semánticas con percentiles utilizados como señal de inspección. Los umbrales definitivos se documentarán tras el análisis exploratorio y se comparará el efecto de filtrado sobre la distribución. El objetivo es evitar que la limpieza elimine precisamente aquellos eventos que explican parte de la heterogeneidad urbana.

\section{Análisis exploratorio previsto}

El análisis exploratorio se estructura en cuatro bloques. El primero cuantifica volumen y cobertura de datos. El segundo estudia distribución temporal. El tercero examina geografía y flujos. El cuarto analiza campos económicos allí donde sean comparables.

\subsection{Volumen y cobertura}

Se reportarán número de ficheros, tamaño, registros y porcentaje de filas válidas por servicio y año. En la versión final, la plataforma habrá procesado \TotalTrips{} viajes y un volumen Bronze de \TotalBronzeGB{}.

\placeholderfigure{Serie mensual del número de viajes por servicio entre 2018 y 2025, generada a partir de \texttt{mart\_mobility\_overview}.}{Evolución mensual del volumen de viajes por tipo de servicio.}

\subsection{Patrones temporales}

Se analizarán perfiles horarios, diferencias laborable/fin de semana, estacionalidad mensual y periodos de ruptura. Es importante representar tanto valores absolutos como índices relativos, dado que la escala entre servicios puede diferir varios órdenes de magnitud.

\subsection{Patrones geográficos}

Los mapas mostrarán densidad de recogidas por zona y cambios relativos entre periodos. También se construirán matrices o rankings de pares origen-destino, evitando visualizaciones que resulten ilegibles por el número de combinaciones.

\subsection{Indicadores económicos}

Para Yellow y Green se estudiarán distribuciones de importe, propina, distancia, duración e ingreso observado por hora ocupada. Estas métricas se utilizarán para el dashboard de oportunidades, pero siempre acompañadas de tamaño muestral para evitar rankings dominados por zonas con pocas observaciones.

\section{Trazabilidad de la fuente al indicador}

Cada tabla Gold conserva las claves necesarias para conocer la procedencia de las agregaciones. En el nivel transaccional se conserva \texttt{source\_file}; en el nivel agregado se mantiene el periodo y servicio. La tabla de control permite reconstruir qué ficheros participaron en una ejecución determinada. Esta trazabilidad es necesaria tanto para depuración como para defensa académica de los resultados.
